\begin{abstract}
\noindent 本研究探讨传染病再生数估计中增加报告病例量的统计收益如何取决于推断目标。在包含环境变异、过度离散后代与二项报告的分支更新过程模型中，由 $m$ 个期望病例构建的条件无偏估计量，对当期瞬时值 $R_t$ 的均方误差为 $A/m$，但对跨期超总体过程均值 $\bar{R}$ 的单期均方误差为 $v_R+A/m$；这是因为 $R_t$ 在时期之间的真实变异构成了推断过程目标时不可消除的跨期抽样变异。未观测的动态报告波动还会引入一项不随 $m$ 下降的方差，即使对当期目标亦存在渐近方差底线；而对有限观察期的实现均值 $\bar{R}_T$，估计误差随总病例量严格衰减为 $A/(Tm)$。利用美国四个流感季早期增长期的住院数据（1748 个州周），仿射损失模型给出的条件截距为正（表观平衡比为每周 37.4 例住院，95\% CI: 19.1, 71.0），但无截距幂律模型在现有观测支持域上的拟合效果相当（$p=0.114$，交叉验证不倾向仿射形式）。受控二项稀释实验验证了计数误差项的 $1/m$ 标度律（斜率增加速率 $c=0.41$，95\% CI: 0.31, 0.53，与泊松基准 $0.45$ 相容；$q\ge0.2$ 时为 $0.48$），且排除了大幅度的截距漂移。州固定效应模型给出的州内平衡比为 17.3（95\% CI: 3.3, 51.4），与横截面比值的差异在统计上无法与抽样误差区分。基于匹配模拟的校准显示，在基准代际间隔 3.6 天代理值下潜在环境方差为 0.040（95\% CI: 0.026, 0.056），在直接测定的 2.2 至 2.9 天下为 0.013 至 0.026，且截距的 41\% 至 51\% 为同周所有州共有。这些诊断验证了计数噪声模型的标度响应，同时表明常规横截面监测数据既无法将渐近方差底线与幂律衰减经验分离，亦无法唯一识别潜在传播环境方差。
\end{abstract}

\begin{keyword}
再生数 \sep 环境随机性 \sep 监测设计 \sep 报告异质性 \sep 流感住院 \sep 推断目标
\end{keyword}
